\documentclass{article}
\usepackage{amsmath, amssymb, amsthm}
\usepackage{hyperref}
\usepackage{geometry}
\geometry{margin=1in}

\title{Erd\H{o}s Problem \#592}
\date{Last updated: 2025-08-31}
\author{Source: erdosproblems.com}

\begin{document}
\maketitle

\noindent\textbf{Status:} OPEN \quad \textbf{Prize: \$1000}

\noindent\textbf{Tags:} set theory, ramsey theory

\medskip

\noindent\textbf{Problem Statement:}

Determine which countable ordinals $\beta$ have the property that, if $\alpha=\omega^{^\beta}$, then in any red/blue colouring of the edges of $K_\alpha$ there is either a red $K_\alpha$ or a blue $K_3$.




Such $\alpha$ (in other words, those such that $\alpha \to (\alpha,3)^2$) are called partition ordinals. {UL} {LI}Specker \cite{Sp57} proved this holds for $\beta=2$ and not for $3\leq \beta &lt;\omega$.{/LI} {LI} Chang \cite{Ch72} proved this holds for $\beta=\omega$.{/LI} {LI} Galvin and Larson \cite{GaLa74} have shown that if $\beta\geq 3$ has this property then $\beta$ must be 'additively indecomposable', so that in particular $\beta=\omega^\gamma$ for some countable ordinal $\gamma$. Galvin and Larson conjecture that every $\beta\geq 3$ of this form has this property.{/LI} {LI}Schipperus \cite{Sc10} have proved this is holds if $\beta=\omega^\gamma$ in which $\gamma$ is a countable ordinal which is the sum of one or two indecomposable ordinals, and this fails to hold if $\gamma$ is the sum of four or more indecomposable ordinals.{/LI} {/UL} The remaining open case appears to be when $\gamma$ is the sum of three indecomposable ordinals. The case $\beta=\omega$ is the subject of [590] , and $\beta=\omega^2$ is the subject of [591] . See also [118] , and [1169] for the case $\alpha=\omega_1^2$.




References

[Ch72] Chang, C. C., A partition theorem for the complete graph on {$\omega\sp{\omega }$} . J. Combinatorial Theory Ser. A (1972), 396-452.

[GaLa74] Galvin, Fred and Larson, Jean, Pinning countable ordinals . Fund. Math. (1974/75), 357-361.

[Sc10] Schipperus, Rene, Countable partition ordinals . Ann. Pure Appl. Logic (2010), 1195-1215.

[Sp57] Specker, Ernst, Teilmengen von Mengen mit Relationen . Comment. Math. Helv. (1957), 302-314.


\bigskip
\noindent\small{Source: \url{https://www.erdosproblems.com/592}}
\end{document}
